    
    

    
    \hypertarget{trotter-systematic-error---hubbard-on-the-square-lattice}{%
\section{Trotter systematic error - Hubbard on the square
lattice}\label{trotter-systematic-error---hubbard-on-the-square-lattice}}

    In this example we use the
\href{https://git.physik.uni-wuerzburg.de/ALF/pyALF}{pyALF} interface to
run ALF with the Mz choice of Hubbard-Stratonovich transformation (i.e.,
coupled to the \(z\)-component of the spin) on a \(6\times 6\) site
square lattice, at \(U/t=4\) half-band filling, and inverse temperature
\(\beta t = 5\).

We carry out a systematic \(\Delta \tau t\) extrapolation keeping
\(\Delta \tau t L_\text{Trotter} = 2\) constant. Recall that the
formulation of the auxiliary field QMC approach is based on the
symmetric Trotter decomposition \[
e^{ -\Delta \tau \left( \hat{A} + \hat{B} \right) }  =  e^{ -\Delta \tau \hat{A}/2}  e^{ -\Delta \tau  \hat{B}  }   e^{ -\Delta \tau \hat{A}/2}  +  \mathcal{O} \left (\Delta  \tau^3\right)
\] The overall error produced by this approximation is of the order
\(\Delta \tau^2\).

Bellow we go through the steps for performing this extrapolation:
setting the simulation parameters, running it and analysing the data. A
reference plot for this analyses is found in
\href{https://git.physik.uni-wuerzburg.de/ALF/ALF/-/blob/master/Documentation/doc.pdf}{ALF
documentation}, Sec. 2.3.2 (Symmetric Trotter decomposition).

\begin{center}\rule{0.5\linewidth}{0.5pt}\end{center}

\textbf{1.} Import \texttt{Simulation} class from the \texttt{py\_alf}
python module, which provides the interface with ALF, as well as
mathematics and plotting packages:

    \begin{tcolorbox}[breakable, size=fbox, boxrule=1pt, pad at break*=1mm,colback=cellbackground, colframe=cellborder]
\prompt{In}{incolor}{1}{\boxspacing}
\begin{Verbatim}[commandchars=\\\{\}]
\PY{k+kn}{from} \PY{n+nn}{py\PYZus{}alf} \PY{k+kn}{import} \PY{n}{Simulation}            \PY{c+c1}{\PYZsh{} Interface with ALF}
\PY{c+c1}{\PYZsh{} }
\PY{k+kn}{import} \PY{n+nn}{numpy} \PY{k}{as} \PY{n+nn}{np}                       \PY{c+c1}{\PYZsh{} Numerical library}
\PY{k+kn}{from} \PY{n+nn}{scipy}\PY{n+nn}{.}\PY{n+nn}{optimize} \PY{k+kn}{import} \PY{n}{curve\PYZus{}fit}     \PY{c+c1}{\PYZsh{} Numerical library}
\PY{k+kn}{import} \PY{n+nn}{matplotlib}\PY{n+nn}{.}\PY{n+nn}{pyplot} \PY{k}{as} \PY{n+nn}{plt}          \PY{c+c1}{\PYZsh{} Plotting library}
\end{Verbatim}
\end{tcolorbox}

    \textbf{2.} Create instances of \texttt{Simulation}, specifying the
necessary parameters, in particular the different \(\Delta \tau\)
values:

    \begin{tcolorbox}[breakable, size=fbox, boxrule=1pt, pad at break*=1mm,colback=cellbackground, colframe=cellborder]
\prompt{In}{incolor}{2}{\boxspacing}
\begin{Verbatim}[commandchars=\\\{\}]
\PY{n}{sims} \PY{o}{=} \PY{p}{[}\PY{p}{]}                                \PY{c+c1}{\PYZsh{} Vector of Simulation instances}
\PY{n+nb}{print}\PY{p}{(}\PY{l+s+s1}{\PYZsq{}}\PY{l+s+s1}{dtau values used:}\PY{l+s+s1}{\PYZsq{}}\PY{p}{)}
\PY{k}{for} \PY{n}{dtau} \PY{o+ow}{in} \PY{p}{[}\PY{l+m+mf}{0.05}\PY{p}{,} \PY{l+m+mf}{0.1}\PY{p}{,} \PY{l+m+mf}{0.15}\PY{p}{]}\PY{p}{:}           \PY{c+c1}{\PYZsh{} Values of dtau}
    \PY{n+nb}{print}\PY{p}{(}\PY{n}{dtau}\PY{p}{)}
    \PY{n}{sim} \PY{o}{=} \PY{n}{Simulation}\PY{p}{(}
        \PY{l+s+s1}{\PYZsq{}}\PY{l+s+s1}{Hubbard}\PY{l+s+s1}{\PYZsq{}}\PY{p}{,}                       \PY{c+c1}{\PYZsh{} Hamiltonian}
        \PY{p}{\PYZob{}}                                \PY{c+c1}{\PYZsh{} Model and simulation parameters for each Simulation instance}
        \PY{l+s+s1}{\PYZsq{}}\PY{l+s+s1}{Model}\PY{l+s+s1}{\PYZsq{}}\PY{p}{:} \PY{l+s+s1}{\PYZsq{}}\PY{l+s+s1}{Hubbard}\PY{l+s+s1}{\PYZsq{}}\PY{p}{,}              \PY{c+c1}{\PYZsh{}    Base model}
        \PY{l+s+s1}{\PYZsq{}}\PY{l+s+s1}{Lattice\PYZus{}type}\PY{l+s+s1}{\PYZsq{}}\PY{p}{:} \PY{l+s+s1}{\PYZsq{}}\PY{l+s+s1}{Square}\PY{l+s+s1}{\PYZsq{}}\PY{p}{,}        \PY{c+c1}{\PYZsh{}    Lattice type}
        \PY{l+s+s1}{\PYZsq{}}\PY{l+s+s1}{L1}\PY{l+s+s1}{\PYZsq{}}\PY{p}{:} \PY{l+m+mi}{6}\PY{p}{,}                         \PY{c+c1}{\PYZsh{}    Lattice length in the first unit vector direction}
        \PY{l+s+s1}{\PYZsq{}}\PY{l+s+s1}{L2}\PY{l+s+s1}{\PYZsq{}}\PY{p}{:} \PY{l+m+mi}{6}\PY{p}{,}                         \PY{c+c1}{\PYZsh{}    Lattice length in the second unit vector direction}
        \PY{l+s+s1}{\PYZsq{}}\PY{l+s+s1}{Checkerboard}\PY{l+s+s1}{\PYZsq{}}\PY{p}{:} \PY{k+kc}{False}\PY{p}{,}           \PY{c+c1}{\PYZsh{}    Whether checkerboard decomposition is used or not}
        \PY{l+s+s1}{\PYZsq{}}\PY{l+s+s1}{Symm}\PY{l+s+s1}{\PYZsq{}}\PY{p}{:} \PY{k+kc}{True}\PY{p}{,}                    \PY{c+c1}{\PYZsh{}    Whether symmetrization takes place}
        \PY{l+s+s1}{\PYZsq{}}\PY{l+s+s1}{ham\PYZus{}T}\PY{l+s+s1}{\PYZsq{}}\PY{p}{:} \PY{l+m+mf}{1.0}\PY{p}{,}                    \PY{c+c1}{\PYZsh{}    Hopping parameter}
        \PY{l+s+s1}{\PYZsq{}}\PY{l+s+s1}{ham\PYZus{}U}\PY{l+s+s1}{\PYZsq{}}\PY{p}{:} \PY{l+m+mf}{4.0}\PY{p}{,}                    \PY{c+c1}{\PYZsh{}    Hubbard interaction}
        \PY{l+s+s1}{\PYZsq{}}\PY{l+s+s1}{ham\PYZus{}Tperp}\PY{l+s+s1}{\PYZsq{}}\PY{p}{:} \PY{l+m+mf}{0.0}\PY{p}{,}                \PY{c+c1}{\PYZsh{}    For bilayer systems}
        \PY{l+s+s1}{\PYZsq{}}\PY{l+s+s1}{beta}\PY{l+s+s1}{\PYZsq{}}\PY{p}{:} \PY{l+m+mf}{5.0}\PY{p}{,}                     \PY{c+c1}{\PYZsh{}    Inverse temperature}
        \PY{l+s+s1}{\PYZsq{}}\PY{l+s+s1}{Ltau}\PY{l+s+s1}{\PYZsq{}}\PY{p}{:} \PY{l+m+mi}{0}\PY{p}{,}                       \PY{c+c1}{\PYZsh{}    \PYZsq{}1\PYZsq{} for time\PYZhy{}displaced Green functions; \PYZsq{}0\PYZsq{} otherwise }
        \PY{l+s+s1}{\PYZsq{}}\PY{l+s+s1}{NSweep}\PY{l+s+s1}{\PYZsq{}}\PY{p}{:} \PY{l+m+mi}{200}\PY{p}{,}                   \PY{c+c1}{\PYZsh{}    Number of sweeps}
        \PY{l+s+s1}{\PYZsq{}}\PY{l+s+s1}{NBin}\PY{l+s+s1}{\PYZsq{}}\PY{p}{:} \PY{l+m+mi}{10}\PY{p}{,}                      \PY{c+c1}{\PYZsh{}    Number of bins}
        \PY{l+s+s1}{\PYZsq{}}\PY{l+s+s1}{Dtau}\PY{l+s+s1}{\PYZsq{}}\PY{p}{:} \PY{n}{dtau}\PY{p}{,}                    \PY{c+c1}{\PYZsh{}    Only dtau varies between simulations, Ltrot=beta/Dtau}
        \PY{l+s+s1}{\PYZsq{}}\PY{l+s+s1}{Mz}\PY{l+s+s1}{\PYZsq{}}\PY{p}{:} \PY{k+kc}{True}\PY{p}{,}                      \PY{c+c1}{\PYZsh{}    If true, sets the M\PYZus{}z\PYZhy{}Hubbard model: Nf=2, N\PYZus{}sum=1,}
        \PY{p}{\PYZcb{}}\PY{p}{,}                               \PY{c+c1}{\PYZsh{}             HS field couples to z\PYZhy{}component of magnetization}
        \PY{n}{alf\PYZus{}dir}\PY{o}{=}\PY{l+s+s1}{\PYZsq{}}\PY{l+s+s1}{\PYZti{}/Programs/ALF}\PY{l+s+s1}{\PYZsq{}}\PY{p}{,}        \PY{c+c1}{\PYZsh{} Local ALF copy, if present}
    \PY{p}{)}
    \PY{n}{sims}\PY{o}{.}\PY{n}{append}\PY{p}{(}\PY{n}{sim}\PY{p}{)}
\end{Verbatim}
\end{tcolorbox}

    \begin{Verbatim}[commandchars=\\\{\}]
dtau values used:
0.05
0.1
0.15
    \end{Verbatim}

    \textbf{3.} Compile ALF, downloading it first if not found locally. This
may take a few minutes:

    \begin{tcolorbox}[breakable, size=fbox, boxrule=1pt, pad at break*=1mm,colback=cellbackground, colframe=cellborder]
\prompt{In}{incolor}{3}{\boxspacing}
\begin{Verbatim}[commandchars=\\\{\}]
\PY{n}{sims}\PY{p}{[}\PY{l+m+mi}{0}\PY{p}{]}\PY{o}{.}\PY{n}{compile}\PY{p}{(}\PY{p}{)}                        \PY{c+c1}{\PYZsh{} Compilation needs to be performed only once}
\end{Verbatim}
\end{tcolorbox}

    \begin{Verbatim}[commandchars=\\\{\}]
Compiling ALF{\ldots} Done.
    \end{Verbatim}

    \textbf{4.} Perform the simulations, as specified in each element of
\texttt{sim}:

    \begin{tcolorbox}[breakable, size=fbox, boxrule=1pt, pad at break*=1mm,colback=cellbackground, colframe=cellborder]
\prompt{In}{incolor}{4}{\boxspacing}
\begin{Verbatim}[commandchars=\\\{\}]
\PY{k}{for} \PY{n}{i}\PY{p}{,} \PY{n}{sim} \PY{o+ow}{in} \PY{n+nb}{enumerate}\PY{p}{(}\PY{n}{sims}\PY{p}{)}\PY{p}{:}
    \PY{n}{sim}\PY{o}{.}\PY{n}{run}\PY{p}{(}\PY{p}{)}                            \PY{c+c1}{\PYZsh{} Perform the actual simulation in ALF}
\end{Verbatim}
\end{tcolorbox}

    \begin{Verbatim}[commandchars=\\\{\}]
Prepare directory "/home/stafusa/ALF/pyALF/Hubbard\_Square\_L1=6\_L2=6\_Checkerboard
=False\_Symm=True\_T=1.0\_U=4.0\_Tperp=0.0\_beta=5.0\_Dtau=0.05\_Mz=True" for Monte
Carlo run.
Resuming previous run.
Run /home/stafusa/Programs/ALF/Prog/Hubbard.out
Prepare directory "/home/stafusa/ALF/pyALF/Hubbard\_Square\_L1=6\_L2=6\_Checkerboard
=False\_Symm=True\_T=1.0\_U=4.0\_Tperp=0.0\_beta=5.0\_Dtau=0.1\_Mz=True" for Monte
Carlo run.
Resuming previous run.
Run /home/stafusa/Programs/ALF/Prog/Hubbard.out
Prepare directory "/home/stafusa/ALF/pyALF/Hubbard\_Square\_L1=6\_L2=6\_Checkerboard
=False\_Symm=True\_T=1.0\_U=4.0\_Tperp=0.0\_beta=5.0\_Dtau=0.15\_Mz=True" for Monte
Carlo run.
Resuming previous run.
Run /home/stafusa/Programs/ALF/Prog/Hubbard.out
    \end{Verbatim}

    \textbf{5.} Calculate the internal energies:

    \begin{tcolorbox}[breakable, size=fbox, boxrule=1pt, pad at break*=1mm,colback=cellbackground, colframe=cellborder]
\prompt{In}{incolor}{5}{\boxspacing}
\begin{Verbatim}[commandchars=\\\{\}]
\PY{n}{ener} \PY{o}{=} \PY{n}{np}\PY{o}{.}\PY{n}{empty}\PY{p}{(}\PY{p}{(}\PY{n+nb}{len}\PY{p}{(}\PY{n}{sims}\PY{p}{)}\PY{p}{,} \PY{l+m+mi}{2}\PY{p}{)}\PY{p}{)}          \PY{c+c1}{\PYZsh{} Matrix for storing energy values}
\PY{n}{dtaus} \PY{o}{=} \PY{n}{np}\PY{o}{.}\PY{n}{empty}\PY{p}{(}\PY{p}{(}\PY{n+nb}{len}\PY{p}{(}\PY{n}{sims}\PY{p}{)}\PY{p}{,}\PY{p}{)}\PY{p}{)}           \PY{c+c1}{\PYZsh{} Matrix for Dtau values, for plotting}
\PY{k}{for} \PY{n}{i}\PY{p}{,} \PY{n}{sim} \PY{o+ow}{in} \PY{n+nb}{enumerate}\PY{p}{(}\PY{n}{sims}\PY{p}{)}\PY{p}{:}
    \PY{n+nb}{print}\PY{p}{(}\PY{n}{sim}\PY{o}{.}\PY{n}{sim\PYZus{}dir}\PY{p}{)}                   \PY{c+c1}{\PYZsh{} Directory containing the simulation output}
    \PY{n}{sim}\PY{o}{.}\PY{n}{analysis}\PY{p}{(}\PY{p}{)}                       \PY{c+c1}{\PYZsh{} Perform default analysis}
    \PY{n}{dtaus}\PY{p}{[}\PY{n}{i}\PY{p}{]} \PY{o}{=} \PY{n}{sim}\PY{o}{.}\PY{n}{sim\PYZus{}dict}\PY{p}{[}\PY{l+s+s1}{\PYZsq{}}\PY{l+s+s1}{Dtau}\PY{l+s+s1}{\PYZsq{}}\PY{p}{]}                             \PY{c+c1}{\PYZsh{} Store Dtau value}
    \PY{n}{ener}\PY{p}{[}\PY{n}{i}\PY{p}{]} \PY{o}{=} \PY{n}{sim}\PY{o}{.}\PY{n}{get\PYZus{}obs}\PY{p}{(}\PY{p}{[}\PY{l+s+s1}{\PYZsq{}}\PY{l+s+s1}{Ener\PYZus{}scalJ}\PY{l+s+s1}{\PYZsq{}}\PY{p}{]}\PY{p}{)}\PY{p}{[}\PY{l+s+s1}{\PYZsq{}}\PY{l+s+s1}{Ener\PYZus{}scalJ}\PY{l+s+s1}{\PYZsq{}}\PY{p}{]}\PY{p}{[}\PY{l+s+s1}{\PYZsq{}}\PY{l+s+s1}{obs}\PY{l+s+s1}{\PYZsq{}}\PY{p}{]}  \PY{c+c1}{\PYZsh{} Store internal energy}
\end{Verbatim}
\end{tcolorbox}

    \begin{Verbatim}[commandchars=\\\{\}]
/home/stafusa/ALF/pyALF/Hubbard\_Square\_L1=6\_L2=6\_Checkerboard=False\_Symm=True\_T=
1.0\_U=4.0\_Tperp=0.0\_beta=5.0\_Dtau=0.05\_Mz=True
Analysing Ener\_scal
Analysing Part\_scal
Analysing Pot\_scal
Analysing Kin\_scal
Analysing Den\_eq
Analysing SpinZ\_eq
Analysing Green\_eq
Analysing SpinXY\_eq
Analysing SpinT\_eq
/home/stafusa/ALF/pyALF/Hubbard\_Square\_L1=6\_L2=6\_Checkerboard=False\_Symm=True\_T=
1.0\_U=4.0\_Tperp=0.0\_beta=5.0\_Dtau=0.1\_Mz=True
Analysing Ener\_scal
Analysing Part\_scal
Analysing Pot\_scal
Analysing Kin\_scal
Analysing Den\_eq
Analysing SpinZ\_eq
Analysing Green\_eq
Analysing SpinXY\_eq
Analysing SpinT\_eq
/home/stafusa/ALF/pyALF/Hubbard\_Square\_L1=6\_L2=6\_Checkerboard=False\_Symm=True\_T=
1.0\_U=4.0\_Tperp=0.0\_beta=5.0\_Dtau=0.15\_Mz=True
Analysing Ener\_scal
Analysing Part\_scal
Analysing Pot\_scal
Analysing Kin\_scal
Analysing Den\_eq
Analysing SpinZ\_eq
Analysing Green\_eq
Analysing SpinXY\_eq
Analysing SpinT\_eq
    \end{Verbatim}

    \begin{tcolorbox}[breakable, size=fbox, boxrule=1pt, pad at break*=1mm,colback=cellbackground, colframe=cellborder]
\prompt{In}{incolor}{6}{\boxspacing}
\begin{Verbatim}[commandchars=\\\{\}]
\PY{n+nb}{print}\PY{p}{(}\PY{l+s+s1}{\PYZsq{}}\PY{l+s+s1}{For Dtau values}\PY{l+s+s1}{\PYZsq{}}\PY{p}{,} \PY{n}{dtaus}\PY{p}{,} \PY{l+s+s1}{\PYZsq{}}\PY{l+s+s1}{the measured energies are:}\PY{l+s+se}{\PYZbs{}n}\PY{l+s+s1}{\PYZsq{}}\PY{p}{,} \PY{n}{ener}\PY{p}{)}
\end{Verbatim}
\end{tcolorbox}

    \begin{Verbatim}[commandchars=\\\{\}]
For Dtau values [0.05 0.1  0.15] the measured energies are:
 [[-29.714802   0.041044]
 [-29.784918   0.043263]
 [-29.818716   0.029992]]
    \end{Verbatim}

    \begin{tcolorbox}[breakable, size=fbox, boxrule=1pt, pad at break*=1mm,colback=cellbackground, colframe=cellborder]
\prompt{In}{incolor}{7}{\boxspacing}
\begin{Verbatim}[commandchars=\\\{\}]
\PY{n}{plt}\PY{o}{.}\PY{n}{errorbar}\PY{p}{(}\PY{n}{dtaus}\PY{o}{*}\PY{o}{*}\PY{l+m+mi}{2}\PY{p}{,} \PY{n}{ener}\PY{p}{[}\PY{p}{:}\PY{p}{,} \PY{l+m+mi}{0}\PY{p}{]}\PY{p}{,} \PY{n}{ener}\PY{p}{[}\PY{p}{:}\PY{p}{,} \PY{l+m+mi}{1}\PY{p}{]}\PY{p}{)}

\PY{k}{def} \PY{n+nf}{func}\PY{p}{(}\PY{n}{x}\PY{p}{,} \PY{n}{y0}\PY{p}{,} \PY{n}{a}\PY{p}{)}\PY{p}{:}
    \PY{k}{return} \PY{n}{y0} \PY{o}{+} \PY{n}{a}\PY{o}{*}\PY{n}{x}\PY{o}{*}\PY{o}{*}\PY{l+m+mi}{2}
\PY{n}{popt1}\PY{p}{,} \PY{n}{pcov} \PY{o}{=} \PY{n}{curve\PYZus{}fit}\PY{p}{(}\PY{n}{func}\PY{p}{,} \PY{n}{dtaus}\PY{p}{,} \PY{n}{ener}\PY{p}{[}\PY{p}{:}\PY{p}{,} \PY{l+m+mi}{0}\PY{p}{]}\PY{p}{,} \PY{n}{sigma}\PY{o}{=}\PY{n}{ener}\PY{p}{[}\PY{p}{:}\PY{p}{,} \PY{l+m+mi}{1}\PY{p}{]}\PY{p}{,} \PY{n}{absolute\PYZus{}sigma}\PY{o}{=}\PY{k+kc}{True}\PY{p}{)}
\PY{n}{perr1} \PY{o}{=} \PY{n}{np}\PY{o}{.}\PY{n}{sqrt}\PY{p}{(}\PY{n}{np}\PY{o}{.}\PY{n}{diag}\PY{p}{(}\PY{n}{pcov}\PY{p}{)}\PY{p}{)}
\PY{n+nb}{print}\PY{p}{(}\PY{n}{popt1}\PY{p}{,} \PY{n}{perr1}\PY{p}{)}
\PY{n}{xs} \PY{o}{=} \PY{n}{np}\PY{o}{.}\PY{n}{linspace}\PY{p}{(}\PY{l+m+mf}{0.}\PY{p}{,} \PY{n}{dtaus}\PY{o}{.}\PY{n}{max}\PY{p}{(}\PY{p}{)}\PY{p}{)}
\PY{n}{plt}\PY{o}{.}\PY{n}{plot}\PY{p}{(}\PY{n}{xs}\PY{o}{*}\PY{o}{*}\PY{l+m+mi}{2}\PY{p}{,} \PY{n}{func}\PY{p}{(}\PY{n}{xs}\PY{p}{,} \PY{o}{*}\PY{n}{popt1}\PY{p}{)}\PY{p}{)}

\PY{n}{plt}\PY{o}{.}\PY{n}{errorbar}\PY{p}{(}\PY{l+m+mi}{0}\PY{p}{,} \PY{n}{popt1}\PY{p}{[}\PY{l+m+mi}{0}\PY{p}{]}\PY{p}{,} \PY{n}{perr1}\PY{p}{[}\PY{l+m+mi}{0}\PY{p}{]}\PY{p}{)}
\end{Verbatim}
\end{tcolorbox}

    \begin{Verbatim}[commandchars=\\\{\}]
[-29.71520822  -4.77614903] [0.04068013 2.44507809]
    \end{Verbatim}

            \begin{tcolorbox}[breakable, size=fbox, boxrule=.5pt, pad at break*=1mm, opacityfill=0]
\prompt{Out}{outcolor}{7}{\boxspacing}
\begin{Verbatim}[commandchars=\\\{\}]
<ErrorbarContainer object of 3 artists>
\end{Verbatim}
\end{tcolorbox}
        
    \begin{center}
    \adjustimage{max size={0.6\linewidth}{0.6\paperheight}}{trotter_error-output_12_2.png}
    \end{center}
    { \hspace*{\fill} \\}
    
    \begin{center}\rule{0.5\linewidth}{0.5pt}\end{center}

\hypertarget{exercises}{%
\subsection{Exercises}\label{exercises}}

\begin{enumerate}
\def\labelenumi{\arabic{enumi}.}
\tightlist
\item
  Try out the four different combinations for \texttt{Checkerboard} and
  \texttt{Symm} settings in order to observe their effect on the output
  and run time. Reference: Sec. 2.3.2 - Symmetric Trotter decomposition
  - of the
  \href{https://git.physik.uni-wuerzburg.de/ALF/ALF/-/blob/master/Documentation/doc.pdf}{ALF
  documentation}, especially Fig. 2.
\end{enumerate}


    % Add a bibliography block to the postdoc
    
    
    
